% Вспомогательные команды и TikZ-стили проекта

% --- tcolorbox ---
\newtcolorbox{newsbox}{
  colback=BPLight,
  colframe=BPLight,
  boxrule=0pt,
  arc=2mm,
  left=2mm,
  right=2mm,
  top=1mm,
  bottom=1mm,
  fontupper=\small
}

% --- Inline-акценты ---
\newcommand{\lead}[1]{\textcolor{BPBlue}{#1}}
\newcommand{\tail}[1]{\textcolor{BPMuted}{#1}}

% --- Общие стили блоков пайплайна ---
\tikzset{
  % Базовый блок (нейтральный)
  pipebox/.style={
    draw=BPGray,
    fill=BPGrayFill,
    rounded corners=2mm,
    minimum width=2.0cm,
    minimum height=1.6cm,
    inner sep=2.5mm,
    align=center,
    font=\scriptsize,
    text=BPText
  },
  % Корпус
  pipecorpus/.style={
    pipebox,
    draw=BPTeal,
    fill=BPTealFill,
    text=BPTeal
  },
  % Сэмплинг — нейтральный, наследует pipebox
  pipesample/.style={
    pipebox
  },
  % Суммаризация
  pipesumm/.style={
    pipebox,
    draw=BPSkyBorder,
    fill=BPSkyFill,
    text=BPText
  },
  % Нейтрализация — главный вклад
  pipeactive/.style={
    pipebox,
    draw=BPAccent,
    fill=BPAccentFill,
    text=BPAccent,
    line width=0.9pt
  },
  % Прогноз / оценка
  pipeforecast/.style={
    pipebox,
    draw=BPNavy,
    fill=BPNavyFill,
    text=BPNavy
  },
  % Пунктирный блок (Источники / Респонденты)
  pipedashed/.style={
    pipebox,
    draw=BPLine,
    fill=BPLight,
    text=BPText,
    dashed
  },
  % Приглушённый блок (неактивный шаг)
  pipemuted/.style={
    pipebox,
    draw=BPMuted!60,
    fill=BPLight,
    text=BPMuted
  },
  % Стрелка между блоками
  pipearrow/.style={
    -{Latex[length=2.2mm,width=1.6mm]},
    draw=BPBlue,
    line width=0.8pt
  },
  % Стрелка приглушённая (вертикальные флоу-диаграммы)
  pipearrowmuted/.style={
    -{Latex[length=2.2mm,width=1.6mm]},
    draw=BPMuted,
    line width=0.8pt
  }
}

% --- Слайд-макет: хедер и футер (plain-фреймы) ---
% Использование: \slideheader{Заголовок слайда}
\newcommand{\slideheader}[1]{%
  \draw[BPBlue, line width=0.33mm] (0,86) -- (960,86);
  \pgfmathsetmacro{\progressx}{48 + 864 * \insertframenumber / \inserttotalframenumber}
  \fill[BPBlue, rounded corners=0.33mm] (48,67) rectangle (\progressx,70);
  \node[anchor=north west, inner sep=0pt, text=BPBlue]
    at (48,28) {\adjustbox{max width=144mm}{\fontsize{17}{21}\selectfont\bfseries #1}};
}

% \slidefooter{Имя автора}{Заголовок презентации}
\newcommand{\slidefooter}[2]{%
  \draw[BPFooter, line width=0.17mm] (0,510) -- (960,510);
  \node[anchor=west, inner sep=0pt, text=BPBlue,
    font=\fontsize{8.25}{10}\selectfont\bfseries]
    at (48,525) {#1};
  \node[anchor=west, inner sep=0pt, text=BPMuted, align=center]
    at (220,525) {\resizebox{112mm}{!}{#2}};
  \node[anchor=east, inner sep=0pt, text=BPMuted,
    font=\fontsize{8.25}{10}\selectfont]
    at (912,525) {\insertframenumber};
}

\setbeamertemplate{frametitle}{%
  \begin{tikzpicture}[remember picture,overlay,x=0.1666667mm,y=-0.1666667mm]
    \begin{scope}[shift={(current page.north west)}]
      \slideheader{\insertframetitle}
    \end{scope}
  \end{tikzpicture}%
}

\setbeamertemplate{footline}{%
  \begin{tikzpicture}[remember picture,overlay,x=0.1666667mm,y=-0.1666667mm]
    \begin{scope}[shift={(current page.north west)}]
      \slidefooter{Сергей Устьянцев}{Извлечение сигнала инфляционных ожиданий из новостного потока с контролем утечки будущих данных}
    \end{scope}
  \end{tikzpicture}%
  \vspace{7mm}%
}
